% Einheit 4 von 4 – Prozentuale Veränderung (Katalog Einheit 5), Kl. 7, schwach
\clearpage
\hypertarget{e4}{}
\einheitenkopf*{Einheit 4 von 4 · Prozentuale Veränderung}
\zweigzeile{Hier lernst du, Preise nach einer Erhöhung oder Senkung auszurechnen und eine Veränderung in Prozent anzugeben · neu oder schon bekannt – je nach Buch · P10 oft · baut auf: Einheit 1 bis 3, Rest zu 100\,\% und Runden (Kennst du schon)}
\setcounter{aufgabe}{38}

\begin{aufgabe}{Ich erkenne, ob „um“ oder „auf“ gemeint ist. „Um 20\,\% gesenkt“ heißt: 20\,\% fallen weg. „Auf 80\,\% gesenkt“ heißt: 80\,\% bleiben. Rechne nicht. Kreuze an, was dasselbe bedeutet.}
\swb{\small grau: es bleiben 80\,\% \\ weiß: 20\,\% fallen weg}{\streifen{80}{$0\,\%$}{$100\,\%$}}
\begin{teile}
\teil Der Preis wird um 40\,\% gesenkt. \\ \kreuz{auf 60\,\% gesenkt}\quad\kreuz{auf 40\,\% gesenkt}
\teil Die Miete steigt auf 110\,\%. \\ \kreuz{um 10\,\% erhöht}\quad\kreuz{um 110\,\% erhöht}
\teil Der Preis wird auf 90\,\% gesenkt. \\ \kreuz{um 90\,\% gesenkt}\quad\kreuz{um 10\,\% gesenkt}
\teil Der Eintritt wird um 25\,\% erhöht. \\ \kreuz{auf 125\,\% erhöht}\quad\kreuz{auf 25\,\% erhöht}
\teil Das Spiel kostet jetzt nur noch 75\,\% vom alten Preis. \\ \kreuz{um 75\,\% gesenkt}\quad\kreuz{um 25\,\% gesenkt}
\end{teile}
\end{aufgabe}

\begin{aufgabe}{Ich erkenne, welcher Wert der alte ist. Der alte Wert ist das Ganze, also 100\,\%. Unterstreiche in jedem Satz den alten Wert.}
\swb{\small alter Wert $= 100\,\%$}{\streifen{100}{$0\,\%$}{$100\,\%$}}
\begin{teile}
\teil Ein Brot kostete 3\,€. Jetzt kostet es 3,30\,€.
\teil Jetzt kostet die Karte 12\,€. Früher waren es 10\,€.
\teil Im Mai hatte der Verein 28 Mitglieder, im Juni sind es 35.
\teil Nach der Senkung kostet die Jacke 45\,€. Vorher kostete sie 60\,€.
\teil Die Miete steigt von 500\,€ auf 540\,€.
\end{teile}
\end{aufgabe}

\begin{aufgabe}{Ich kann ausrechnen, wie viel etwas nach einer Erhöhung oder Senkung kostet. Rechne zuerst aus, wie viel Euro dazukommen. Zähle sie dann zum alten Preis dazu.}
\swb{\small 20\,€ um 10\,\% erhöht\\[2pt]\beispiel{10\,\% \text{ von } 20\,€ &= 2\,€ \\ 20\,€ + 2\,€ &= 22\,€}}{\streifenfrage{10}{$0\,€$}{$20\,€$}}
\anw{Der Preis wird immer um 10\,\% erhöht.}
\swz{alter Preis 30\,€}{\streifenfrage{10}{$0\,€$}{$30\,€$}}
\swz{alter Preis 50\,€}{\streifenfrage{10}{$0\,€$}{$50\,€$}}
\swz{alter Preis 70\,€}{\streifenfrage{10}{$0\,€$}{$70\,€$}}
\swz{alter Preis 90\,€}{\streifenfrage{10}{$0\,€$}{$90\,€$}}
\swfrage{Was bleibt gleich, was ändert sich?}
\end{aufgabe}

\begin{aufgabe}{Ich kann ausrechnen, wie viel etwas nach einer Erhöhung oder Senkung kostet – weiter. Bei einer Senkung ziehst du den Prozentwert vom alten Preis ab.}
\swz{80\,€ werden um 25\,\% gesenkt.}{\streifen[20]{75}{$0\,€$}{$80\,€$}}
\anw{Bei „auf“ rechnest du den neuen Preis gleich aus.}
\swz{70\,€ werden auf 90\,\% gesenkt.}{\streifenfrage{90}{$0\,€$}{$70\,€$}}
\end{aufgabe}

\begin{aufgabe}{Ich kann den neuen Preis in einem Schritt mit einer Kommazahl ausrechnen (kommt in Klasse 9). Rechne aus, wie viel Prozent vom alten Preis der neue Preis ist. Schreibe das als Kommazahl und nimm sie mal den alten Preis.}
\swb{\small 40\,€ um 10\,\% erhöht\\[2pt]\beispiel{100\,\% + 10\,\% &= 110\,\% = 1{,}1 \\ 1{,}1 \cdot 40\,€ &= 44\,€}}{\zahlenstrahl[xmin=0.8,xmax=1.35,xstep=0.05,karo=0.8]{1/100\,\%, 1.1/110\,\%}}
\swz{70\,€ werden um 30\,\% erhöht.}{\zahlenstrahl[xmin=0.8,xmax=1.35,xstep=0.05,karo=0.8]{1/100\,\%, 1.3/130\,\%}}
\swz{36\,€ werden um 25\,\% gesenkt.}{\zahlenstrahl[xmin=0,xmax=1.3,xstep=0.25,karo=1.6]{0.75/75\,\%, 1/100\,\%}}
\end{aufgabe}

\begin{aufgabe}{Ich kann ausrechnen, um wie viel Prozent sich ein Preis verändert hat. Rechne zuerst den Unterschied aus. Teile ihn dann durch den alten Preis.}
\swb{\small von 50\,€ auf 60\,€\\[2pt]\beispiel{60\,€ - 50\,€ &= 10\,€ \\ 10 : 50 &= 0{,}2 = 20\,\%}}{\streifenfrage{20}{$0\,€$}{$50\,€$}}
\swz{von 40\,€ auf 50\,€}{\streifenfrage[20]{25}{$0\,€$}{$40\,€$}}
\swz{von 80\,€ auf 68\,€}{\streifenfrage[20]{15}{$0\,€$}{$80\,€$}}
\end{aufgabe}

\begin{aufgabe}{Ich kann ausrechnen, wie viel etwas vor einer Änderung gekostet hat. Überlege zuerst: Wie viel Prozent vom alten Preis ist der neue Preis?}
\swz{Nach 25\,\% Rabatt kostet ein Spiel 30\,€. Wie viel kostete es vorher?}{\streifen[20]{75}{$0\,€$}{$?$}}
\swz{Nach einer Erhöhung um 10\,\% kostet eine Karte 5,50\,€. Wie viel kostete sie vorher?}{\zahlenstrahl[xmin=0.8,xmax=1.35,xstep=0.05,karo=0.8]{1/100\,\%, 1.1/110\,\%}}
\end{aufgabe}

\begin{aufgabe}{Ich kann Preise mit und ohne Mehrwertsteuer ausrechnen (neu in diesem Jahr). Der Preis ohne Steuer heißt Nettopreis, er ist 100\,\%. Der Preis mit Steuer heißt Bruttopreis.}
\swz{Ein Regal kostet netto 200\,€. Dazu kommen 19\,\% Mehrwertsteuer. Wie viel kostet es brutto?}{\zahlenstrahl[xmin=0.9,xmax=1.275,xstep=0.05,karo=1.3]{1/100\,\%, 1.19/119\,\%}}
\swz{Ein Buch kostet brutto 10,70\,€ mit 7\,\% Mehrwertsteuer. Wie viel kostet es netto?}{\zahlenstrahl[xmin=0.9,xmax=1.275,xstep=0.05,karo=1.3]{1/100\,\%, 1.07/107\,\%}}
\end{aufgabe}

\begin{aufgabe}{Ich kann zwei Prozentangaben richtig vergleichen. Bei einer Schulwahl bekam eine Liste erst 40\,\% der Stimmen, bei der nächsten Wahl 46\,\%. Den Unterschied zweier Prozentangaben nennt man Prozentpunkte.}
\swz{Um wie viele Prozentpunkte ist der Anteil gestiegen?}{\streifen{40}{$0\,\%$}{$100\,\%$}\par\smallskip\streifen[0]{46}{$0\,\%$}{$100\,\%$}}
\swz{Um wie viel Prozent ist der Anteil gestiegen? Der alte Anteil 40\,\% ist das Ganze.}{\streifenfrage[20]{15}{$0$}{$40\,\%$}}
\end{aufgabe}

\begin{aufgabe}{Ich kann eine Steigung in Prozent verstehen und ausrechnen (kommt in Klasse 10). Die Steigung in Prozent sagt, wie viele Meter es auf 100\,m waagerechter Strecke nach oben geht. Die Skizze ist nicht maßstabsgetreu.}
\swz{Eine Straße hat 5\,\% Steigung. Auf 100\,m waagerechter Strecke geht es \leerfeld[m] nach oben.}{\dreieckrw{5}{0.8}{100\,\text{m}}{?}{}}
\swz{Eine Straße steigt auf 400\,m waagerechter Strecke um 30\,m an. Wie viel Prozent Steigung hat sie?}{\dreieckrw{5}{0.8}{400\,\text{m}}{30\,\text{m}}{}}
\end{aufgabe}

\begin{aufgabe}{Ich kann ausrechnen, um wie viel Prozent sich ein Preis verändert hat – weiter.}
\begin{teile}
\teil Ein Liter Milch kostete im Januar 1,15\,€. Im März kostete er 1,24\,€. Um wie viel Prozent ist der Preis gestiegen? Runde auf eine Stelle nach dem Komma. (P10 2026)
\end{teile}
\end{aufgabe}

\begin{aufgabe}{Ich finde den Fehler bei einer Veränderung in Prozent. Ein Preis sinkt von 50\,€ auf 40\,€. Leo rechnet so:}
\rechnung{50\,€ - 40\,€ &= 10\,€ \\ 10 : 40 &= 0{,}25 = 25\,\%}
\swz{Finde den Fehler, benenne ihn und korrigiere die Rechnung.}{\streifenfrage{20}{$0\,€$}{$50\,€$}}
\swz{Eine Fahrkarte kostete 2,00\,€. Jetzt kostet sie 2,30\,€. Um wie viel Prozent ist sie teurer geworden?}{\streifenfrage[20]{15}{$0\,€$}{$2{,}00\,€$}}
\end{aufgabe}

\begin{aufgabe}{Ich kann begründen, welcher Wert 100\,\% ist.}
\swz{Ein Preis steigt von 20\,€ auf 25\,€. Warum teilt man den Unterschied durch 20 und nicht durch 25? Begründe.}{\streifen[4]{25}{$0\,€$}{$20\,€$}}
\swz{Ein Preis steigt von 80\,€ auf 100\,€ und fällt danach wieder auf 80\,€. Anna sagt: „Er ist um 25\,\% gestiegen und danach um 25\,\% gefallen.“ Stimmt das? Begründe.}{\streifen[20]{25}{$0\,€$}{$80\,€$}\par\smallskip\streifen{20}{$0\,€$}{$100\,€$}}
\end{aufgabe}

\begin{aufgabe}{Ich kann entscheiden, ob sich zwei Änderungen aufheben. Ein Laden erhöht alle Preise um 10\,\%. Einen Monat später gibt er 10\,\% Rabatt auf alle Preise. Eine Lampe kostete am Anfang 100\,€.}
\swz{Wie viel kostet die Lampe nach der Erhöhung?}{\streifenfrage{10}{$0\,€$}{$100\,€$}}
\swz{Wie viel kostet die Lampe nach dem Rabatt?}{\streifenfrage{10}{$0\,€$}{$?$}}
\swz{Kostet die Lampe jetzt wieder so viel wie am Anfang? Begründe.}{}
\end{aufgabe}

\uebersichtskasten{Erhöhung um p\,\%: neuer Wert $=$ alter Wert $\cdot\,(1 + \tfrac{p}{100})$; Senkung: $\cdot\,(1 - \tfrac{p}{100})$. \ 250\,€ um 12\,\% erhöht: $250 \cdot 1{,}12 = 280$\,€; \ um 12\,\% gesenkt: $250 \cdot 0{,}88 = 220$\,€ \\ Veränderung in Prozent: (neu $-$ alt) : alt. \quad von 60 auf 69: $9 : 60 = 0{,}15 = 15\,\%$ \\ Prozentpunkte: Unterschied zweier Prozentangaben. \quad von 30\,\% auf 33\,\%: 3 Prozentpunkte (das sind 10\,\% mehr) \\ Steigung in Prozent: Höhe : waagerechte Strecke. \quad 8\,m auf 100\,m $=$ 8\,\%}
